% common/pool/立体几何/难题/q_solid_rhombus_expand_volume.tex
% ★★★ 难题

\newcommand{\qSolidRhombusExpandVolumeStem}{
已知菱形 $ABCD$ 的边长为 $2$，$\angle ABC=30^\circ$。点 $P$ 在菱形 $ABCD$ 的内部及边界上运动，空间中的点 $Q$ 满足 $|PQ|\leqslant 1$。求点 $Q$ 轨迹所围成的几何体的体积。
}

\newcommand{\qSolidRhombusExpandVolumeOptions}{
A. $\frac{14}{3}\pi+4\sqrt{3}$ \quad
B. $\frac{14}{3}\pi+4$ \quad
C. $\frac{16}{3}\pi+4\sqrt{3}$ \quad
D. $\frac{16}{3}\pi+4$
}

\newcommand{\qSolidRhombusExpandVolumeFigure}[1][1]{
\begin{tikzpicture}[scale=#1, line join=round, line cap=round, >=stealth]
    % 定义菱形 ABCD 顶点
    \coordinate (B) at (0, 0);
    \coordinate (C) at (2, 0);
    \coordinate (A) at ({2*cos(30)}, {2*sin(30)});
    \coordinate (D) at ({2+2*cos(30)}, {2*sin(30)});
    
    % 绘制边框
    \draw[thick] (A) -- (B) -- (C) -- (D) -- cycle;
    
    % 标记顶点
    \node[above left] at (A) {$A$};
    \node[below left] at (B) {$B$};
    \node[below right] at (C) {$C$};
    \node[above right] at (D) {$D$};
    
    % 标出角度和边长
    \draw (0.4,0) arc (0:30:0.4);
    \node at (0.7, 0.15) {\footnotesize $30^\circ$};
\end{tikzpicture}
}

\newcommand{\qSolidRhombusExpandVolumeAnswer}{
D
}

\newcommand{\qSolidRhombusExpandVolumeSolution}{
\textbf{分析外扩几何体的构成：} \\
根据题意，点 $P$ 在菱形 $ABCD$ 及其内部运动，点 $Q$ 满足 $|PQ| \leqslant 1$，即点 $Q$ 位于以 $P$ 为球心、$1$ 为半径的球体内部或边界上。由于 $P$ 遍历整个菱形区域，点 $Q$ 的轨迹围成的几何体可以看作是菱形区域在空间中向各个方向“外扩”距离 $1$ 得到的图形。它可拆分为三部分：
\begin{enumerate}
    \item \textbf{上下平移部分（柱体）}：\\
    当 $P$ 在菱形 $ABCD$ 内部运动时，$Q$ 在垂直于菱形平面的上下两侧形成一个高为 $2$（上下各 $1$）、底面为菱形的直四棱柱。\\
    其体积为 $V_1 = S_{\text{菱形}} \times h = (2 \times 2 \times \sin 30^\circ) \times 2 = 2 \times 2 = 4$。
    
    \item \textbf{四边外扩部分（半圆柱）}：\\
    当 $P$ 在菱形的四条边上运动时，$Q$ 形成四个半圆柱，其半径为 $1$，高分别为菱形的四条边长（均为 $2$）。这四个半圆柱可以拼成两个完整的圆柱。\\
    总体积为 $V_2 = 2 \times (\pi \times 1^2 \times 2) = 4\pi$。
    
    \item \textbf{四角外扩部分（球块）}：\\
    当 $P$ 在菱形的四个顶点处时，$Q$ 形成四个球形凸角（每个均为对应顶点处球面的一部分）。根据凸多边形外角和为 $360^\circ$ 的性质，这四个顶角处外扩的球块拼接起来，恰好是一个完整的半径为 $1$ 的球的体积。\\
    体积为 $V_3 = \frac{4}{3}\pi \times 1^3 = \frac{4}{3}\pi$。
\end{enumerate}

\textbf{总体积计算：} \\
综上所述，该几何体的总体积为：
\[
V = V_1 + V_2 + V_3 = 4 + 4\pi + \frac{4}{3}\pi = \frac{16}{3}\pi + 4
\]
故选 D。
}

\newcommand{\qSolidRhombusExpandVolumeDiagnosis}{
\textbf{学生可能存在的问题：}
1. \textbf{空间想象能力不足}：无法把“动点距离”转化为“区域外扩”，误以为只求一个球体或者圆柱体，漏掉柱体和平移部分。
2. \textbf{角落处理错误}：忽略了菱形顶点处的外扩部分是一个完整的球（拼接后），错算成四个半球或其他比例。
3. \textbf{菱形面积计算错误}：未正确使用正弦定理公式 $S = ab\sin C$，误将边长直接相乘或将高当成边长导致算错底面积。
}

\newcommand{\qSolidRhombusExpandVolumeTeachingNotes}{
\textbf{课堂教学提示：}
1. \textbf{降维类比建模型}：引导学生先回顾二维平面中“线段外扩距离 $d$”的轨迹（跑道形状：一个矩形加两个半圆），进而推广到三维空间中的“面外扩”。
2. \textbf{三部分拆解法}：通过画草图或用手势比划，将几何体严格拆解为“中柱”、“四边半圆柱”、“四角球块”。
3. \textbf{外角和定理巧解}：特别强调，无论底面是什么凸多边形（三角形、矩形、五边形等），其顶点处的球块拼接起来必定是一个完整的球。可以作为“距离轨迹专题”中的通法结论让学生记住。
}
